\chapter{Introducción al mundo}

\section*{Introducción para los jugadores}

\begin{quote}
\textit{El aire en el exoplaneta Ur-Bare siempre sabe a metal y a sangre. Fuera, las tormentas de polvo rojo devoran cualquier cosa que no esté forjada en titanio. Dentro, bajo la inmensa cúpula que protege nuestra colonia, la lluvia jamás descansa. Es un sirimiri sintético y ácido que oxida los pulmones y el alma. Bienvenidos a NeoBilbao.}
\end{quote}

Los jugadores son los\emph{ Langileak}, obreros que viven en los Zócalos Industriales, bajo la cúpula dorada de la sociedad de la elite. Trabajan jornadas infrahumanas en la Ría de Acero para alimentar los altos hornos y mantener encendidas las luces de la Cúpula Dorada, una ciudad de titanio y privilegio que nunca verán de cerca. El motivo es simple y brutal: el Etxearen Zina, el Juramento del Linaje. Sus hijos han sido llevados a la ciudad superior; allí tienen comida, educación y un futuro seguro. A cambio, los padres pagan con su cuerpo, su trabajo y su esperanza.

La verdad que ignoran, sin embargo, es mucho peor: el implante neuronal que llevan en la nuca, el\emph{ Velo}, no es un simple aparato de comunicación. Es una herramienta de control, una pantalla de dopamina que proyecta a los padres hábiles simulaciones hiperrealistas de sus hijos felices. Las voces, los mensajes holográficos y los recordatorios diario no provienen de unos niños sanos. Provienen de una máquina que se alimenta del dolor colectivo y de la culpabilidad de quienes siguen trabajando por ellos.

\subsection*{La ambición de la colonia}

La razón por la que la humanidad levantó este infierno en un mundo remoto fue el Burdin-Urdin, el Acero Azul. Un metal superconductor raro y vital para la tecnología interestelar y la supervivencia de la civilización. Bajo los Zócalos, enterrado en la garganta tectónica, yace el núcleo de su valor. Sin ese metal, la tecnología terrestre y la infraestructura orbital se apagarían. Para los trabajadores, sin embargo, este metal no es progreso. Es una condena.

\subsection*{El prólogo audiovisual}

El siguiente guion puede usarse como prólogo inmersivo antes del Capítulo 1. Sirve para presentar el tono, la opresión industrial y la caída de la mentira.

\begin{CajaCinematica}
\textbf{Título:} NeoBilbao: El despertar\par
\textbf{Duración aproximada:} 1 minuto 20 segundos\par
\textbf{Tono visual:} Ciberpunk industrial, lluvia ácida, humo y metal candente.

\textbf{Escena 1: La Ría de Acero (0:00-0:20)}\par
Visual: toma lenta y descendente sobre el cañón industrial. El río de metal fundido burbujea en el fondo. Sobre él, plataformas de carga se deslizan entre la niebla. La Cúpula Dorada aparece a lo lejos como un espejismo de titanio.

\textbf{Voz en off:}\par
\textit{«El aire en el exoplaneta Ur-Bare siempre sabe a metal y a sangre. Fuera, las tormentas de polvo rojo devoran cualquier cosa que no esté forjada en titanio. Pero aquí abajo, bajo el cielo de acero de los Zócalos Industriales, la lluvia ácida nunca cesa... Bienvenidos a NeoBilbao.»}

\textbf{Escena 2: El Juramento (0:20-0:40)}\par
Obreros con trajes de presión y respiradores, avanzando por pasarelas. El implante del Velo brilla en sus nucas. Un escáner Tártalo sobrevuele la fila con su haz rojo.

\textbf{Voz en off:}\par
\textit{«Nacemos, trabajamos y morimos en la oscuridad de la fundición para que Orekak-Corp mantenga las luces de un paraíso que nunca pisaremos. ¿Por qué soportamos este infierno? Por el Etxearen Zina... El Juramento del Linaje. Por nuestros hijos, reubicados en el lujo y la seguridad de la Cúpula Dorada. Pagamos su estancia con cada hora de nuestras vidas en esta fosa.»}

\textbf{Escena 3: El cordón umbilical (0:40-0:55)}\par
Primer plano del visor del obrero. Le llega un holograma del hijo sonriendo. El efecto dopaminérgico del Velo relaja sus hombros.

\textbf{Voz en off:}\par
\textit{«El Velo es nuestro único consuelo. El cordón umbilical que nos une a ellos. Escucharles darnos las gracias a través de los hologramas en cada descanso es lo único que nos mantiene cuerdos... Lo único por lo que vale la pena este aire envenenado.»}

\textbf{Escena 4: El cortocircuito (0:55-1:10)}\par
Un fogonazo de plasma sacude la fundición. El holograma del Velo se rompe en estática, líneas de código y registros de error.

\textbf{Voz en off:}\par
\textit{«Pero hoy, el sistema ha fallado... Un chispazo. El dolor... y luego la nada. No hay vídeo. No hay señal. Solo estática.»}

\textbf{Escena 5: El despertar (1:10-1:20)}\par
Texto sobre negro, con la interferencia de una radio antigua:

\textbf{NEOBILBAO: EL JURAMENTO DEL LINAJE\par Capítulo 1: El despertar}

\textit{«Aita... Ama... despertad...».}
\end{CajaCinematica}

\section*{Resumen rápido para el Director}

No desveles la verdad del Velo ni la cosecha pediátrica antes de la primera sesión. Mantén el peso de la mentira durante la primera parte de la campaña. Cuando los personajes vean la estática, la corrupción del código y el horror del núcleo, la campaña se convertirá en una odisea de denuncia, rabia y supervivencia.
